% ============================================================================
% Fichier  : partie2/P02_C09_quantitatif.tex
% Titre    : L'Analyse Quantitative : FAIR et Pertes Annuelles
% Auteur   : MINKA MI NGUIDJOÏ Thierry Emmanuel
% Version  : 1.0 -- novembre 2026
% ============================================================================
% Sources : Freund et Jones (2014), FAIR ; Hubbard et Seiersen (2016).
% Estimations de LiZbiZ : cas fictif ; moyennes PERT calculées et vérifiées.
% ============================================================================

\chapter{L'Analyse Quantitative : FAIR et Pertes Annuelles}
\label{chap:quantitatif}
\chaptermark{FAIR et pertes annuelles}

\epigraph{%
    « Si vous savez qu'un chiffre est entre deux bornes,\\
    vous savez déjà quelque chose de mesurable. »%
}{Principe du cours, d'après D.~Hubbard}

\niveauChapitre{Intermédiaire à expert}{Chapitres~\ref{chap:risque},
\ref{chap:perception} et~\ref{chap:matrices}.}

\begin{boxobjectifs}
\begin{itemize}
    \item Décomposer un risque selon la taxonomie FAIR (fréquence et
          ampleur des pertes).
    \item Estimer chaque facteur par un minimum, une valeur la plus
          probable et un maximum.
    \item Calculer une moyenne PERT et une perte annuelle attendue.
    \item Comparer deux options de traitement en pertes évitées.
    \item Présenter une estimation avec ses hypothèses et ses limites.
\end{itemize}
\end{boxobjectifs}

\bigskip

% ============================================================================
\section*{Question directrice : combien le risque de rançongiciel coûtait-il à LiZbiZ chaque année~?}
% ============================================================================

En juin, le conseil demande un chiffre~: « combien ce risque nous
coûte-t-il~? » Une cote « 12, intolérable » ne répond pas. Mme~MOUAFO
applique la méthode FAIR au risque R-02, avec les directeurs et le
prestataire, sur deux séances.

\medskip
\textbf{Plan du chapitre.}
\begin{enumerate}
    \item La taxonomie FAIR (section~\ref{sec:c09-taxonomie}).
    \item Estimer par trois points (section~\ref{sec:c09-estimer}).
    \item La fréquence des pertes (section~\ref{sec:c09-frequence}).
    \item L'ampleur des pertes (section~\ref{sec:c09-ampleur}).
    \item La perte annuelle attendue (section~\ref{sec:c09-ale}).
    \item Comparer des options (section~\ref{sec:c09-options}).
\end{enumerate}

\separateur

% ============================================================================
\section{La Taxonomie FAIR}
\label{sec:c09-taxonomie}
% ============================================================================

FAIR (\emph{Factor Analysis of Information Risk}) décompose le risque en
facteurs que l'on peut estimer séparément~\cite{freund2014fair}.

\begin{center}
\begin{tikzpicture}[
    b/.style = {draw=CouleurBordDef, fill=CouleurDef, rounded corners=3pt,
        minimum width=2.6cm, minimum height=0.7cm, align=center,
        font=\footnotesize\color{CouleurPrimaire}},
    arrow/.style = {thick, color=CouleurSecondaire}]
\node[b, fill=CouleurPrimaire, text=white] (r) at (0, 0) {\color{white}\textbf{Risque}};
\node[b] (lef) at (-3.4, -1.3) {Fréquence des\\pertes (LEF)};
\node[b] (lm) at (3.4, -1.3) {Ampleur des\\pertes (LM)};
\node[b] (tef) at (-5.0, -2.8) {Fréquence des\\menaces (TEF)};
\node[b] (vul) at (-1.8, -2.8) {Vulnérabilité};
\node[b] (pl) at (1.8, -2.8) {Pertes\\primaires};
\node[b] (sl) at (5.0, -2.8) {Pertes\\secondaires};
\draw[arrow] (r) -- (lef); \draw[arrow] (r) -- (lm);
\draw[arrow] (lef) -- (tef); \draw[arrow] (lef) -- (vul);
\draw[arrow] (lm) -- (pl); \draw[arrow] (lm) -- (sl);
\end{tikzpicture}
\end{center}
\vspace{-0.4em}
\begin{center}{\small\itshape Figure~9.1. La taxonomie FAIR (simplifiée)}\end{center}

\begin{boxdef}
\textbf{Facteurs FAIR.} La \termecle{fréquence des pertes} (LEF) est le
nombre d'événements de perte par an~; elle est le produit de la
\emph{fréquence des menaces} (TEF, tentatives par an) et de la
\emph{vulnérabilité} (probabilité qu'une tentative réussisse). L'\termecle{ampleur
des pertes} (LM) réunit les pertes \emph{primaires} (subies par
l'organisation) et \emph{secondaires} (dues aux réactions des tiers).
\end{boxdef}

\begin{tableaucours}{L{3.4cm}L{8.8cm}}{%
    \tableauHeader{Forme de perte} &
    \tableauHeader{Exemple chez LiZbiZ}
}
Productivité & Marge perdue pendant l'arrêt, heures supplémentaires \\
Réponse & Prestataires de remise en état, conseil, communication \\
Remplacement & Matériel, licences \\
Amendes et jugements & Condamnations, sanctions \\
Avantage concurrentiel & Tarifs ou secrets divulgués \\
Réputation & Clients perdus, contrats non renouvelés \\
\end{tableaucours}
\vspace{-0.8em}
\begin{center}{\small\itshape Tableau~9.1. Les six formes de perte de FAIR}\end{center}

% ============================================================================
\section{Estimer par Trois Points}
\label{sec:c09-estimer}
% ============================================================================

Chaque facteur est estimé par un \emph{minimum}, une \emph{valeur la plus
probable} et un \emph{maximum}, par des experts calibrés
(Chapitre~\ref{chap:perception}). La loi \termecle{PERT} en tire une moyenne~:
\[
\text{moyenne} = \frac{\text{min} + 4 \times \text{plus probable} + \text{max}}{6}.
\]
Elle pondère la valeur la plus probable sans ignorer les extrêmes.

\begin{boxalert}
\textbf{Violation typique.} Demander à l'expert « le » chiffre.

\smallskip
\textbf{Règle~:} demander trois valeurs et les raisons des bornes. Un
chiffre unique cache l'incertitude~; trois valeurs la montrent, et la
simulation (Chapitre~\ref{chap:simulation}) la propage.
\end{boxalert}

% ============================================================================
\section{La Fréquence des Pertes}
\label{sec:c09-frequence}
% ============================================================================

\begin{tableaucours}{L{4.4cm}cccc}{%
    \tableauHeader{Facteur} &
    \tableauHeader{Min.} &
    \tableauHeader{Plus prob.} &
    \tableauHeader{Max.} &
    \tableauHeader{Moyenne}
}
TEF (tentatives sérieuses par an) & 2 & 2,5 & 4 & 2,67 \\
Vulnérabilité sans second facteur & 5~\% & 10~\% & 20~\% & 10,8~\% \\
Vulnérabilité avec second facteur & 1~\% & 2~\% & 5~\% & 2,3~\% \\
\end{tableaucours}
\vspace{-0.8em}
\begin{center}{\small\itshape Tableau~9.2. Fréquence des pertes de R-02 (estimations de juin)}\end{center}

\textbf{Calcul.} LEF sans second facteur~: $2{,}67 \times 0{,}108 \approx
0{,}29$ événement par an, soit un tous les trois à quatre ans environ.
Avec second facteur~: $2{,}67 \times 0{,}023 \approx 0{,}06$, un tous les
seize ans environ. Les valeurs les plus probables donnent $2{,}5 \times
0{,}10 = 0{,}25$, l'hypothèse du Chapitre~\ref{chap:risque}.

% ============================================================================
\section{L'Ampleur des Pertes}
\label{sec:c09-ampleur}
% ============================================================================

\begin{tableaucours}{L{4.4cm}cccc}{%
    \tableauHeader{Perte (millions de FCFA)} &
    \tableauHeader{Min.} &
    \tableauHeader{Plus prob.} &
    \tableauHeader{Max.} &
    \tableauHeader{Moyenne}
}
\multicolumn{5}{l}{\textit{Pertes primaires}} \\
Productivité & 30 & 60 & 150 & 70,0 \\
Réponse & 40 & 80 & 120 & 80,0 \\
Remplacement & 10 & 16 & 40 & 19,0 \\
\multicolumn{5}{l}{\textit{Pertes secondaires (probabilité 0,8 par événement)}} \\
Pénalités contractuelles & 20 & 74 & 120 & 72,7 \\
Réputation (clients perdus) & 0 & 110 & 250 & 115,0 \\
Amendes et jugements & 0 & 0 & 50 & 8,3 \\
\end{tableaucours}
\vspace{-0.8em}
\begin{center}{\small\itshape Tableau~9.3. Ampleur des pertes de R-02}\end{center}

\textbf{Calcul.} Pertes primaires moyennes~: $70 + 80 + 19 = 169$~millions.
Pertes secondaires moyennes~: $72{,}7 + 115 + 8{,}3 = 196$~millions, à
pondérer par leur probabilité~: $0{,}8 \times 196 \approx 157$. Ampleur
moyenne par événement~: $169 + 157 \approx 326$~millions de FCFA.

\begin{boxscenario}{Opération GAROUA. L'estimation confrontée au réel}
\textbf{Constatation.} Coût estimé de la crise de mars~: 312~millions
(Tableau~0.4), dans l'intervalle de l'estimation FAIR et proche de sa
moyenne (326~millions).

\medskip
\textbf{Analyse} (confiance moyenne). Un seul événement ne valide pas un
modèle~; il montre seulement que l'ordre de grandeur n'était pas absurde.
Les experts ont d'ailleurs utilisé le coût de mars pour estimer certaines
valeurs~: la concordance est en partie construite.
\end{boxscenario}

% ============================================================================
\section{La Perte Annuelle Attendue}
\label{sec:c09-ale}
% ============================================================================

\begin{boxdef}
\termecle{Perte annuelle attendue} (FAIR)~: produit de la fréquence
moyenne des pertes et de l'ampleur moyenne d'un événement. Elle compare
des risques et des options~; elle ne dit rien de la pire année, que donne
la simulation.
\end{boxdef}

\begin{tableaucours}{L{4.6cm}R{2.4cm}R{2.4cm}R{2.4cm}}{%
    \tableauHeader{} &
    \tableauHeader{LEF (par an)} &
    \tableauHeader{LM (millions)} &
    \tableauHeader{Perte annuelle}
}
Sans second facteur & 0,289 & 326 & 94 \\
Avec second facteur & 0,062 & 326 & 20 \\
\end{tableaucours}
\vspace{-0.8em}
\begin{center}{\small\itshape Tableau~9.4. Perte annuelle attendue de R-02 (millions de FCFA)}\end{center}

% ============================================================================
\section{Comparer des Options}
\label{sec:c09-options}
% ============================================================================

Le second facteur coûte 6~millions par an. Il évite environ
$94 - 20 = 74$~millions de pertes annuelles attendues. Le rapport entre
pertes évitées nettes et coût~:
\[
\frac{74 - 6}{6} \approx 11.
\]
Chaque franc dépensé en évite, en moyenne, une dizaine. Même avec les
bornes les plus défavorables (TEF de 2, vulnérabilité de 5~\% sans et de
5~\% avec, ce qui annule l'effet), le calcul montre de quoi dépend la
décision~: la réduction de la vulnérabilité, qu'il faut donc vérifier.

\begin{boxobjet}{R-02 · Estimation FAIR (juin 2026)}
\textbf{Scénario~:} rançongiciel par l'accès distant du prestataire.
\textbf{Estimation~:} 0,29 événement par an, 326~millions par événement,
94~millions de perte annuelle attendue. \textbf{Traitement~:} second
facteur (perte attendue ramenée à 20~millions). \textbf{Propriétaire~:}
M.~NGUELE. \textbf{Indicateur~:} part des connexions distantes avec second
facteur (cible~: 100~\%).
\end{boxobjet}

\begin{boxmetier}{Gestionnaire de risques}
\textbf{Montrer les hypothèses sur la même page que le chiffre.} Mme~MOUAFO
présente au conseil la perte annuelle attendue, l'intervalle des
estimations et les deux hypothèses dont le résultat dépend le plus.
\end{boxmetier}

\separateur

\begin{boxresume}
\textbf{Ce qu'il faut retenir de ce chapitre~:}
\begin{itemize}
    \item \textbf{FAIR}~: risque = fréquence des pertes × ampleur~;
          fréquence = menaces × vulnérabilité.
    \item \textbf{Six formes de perte}, primaires et secondaires.
    \item \textbf{Trois points} par facteur, moyenne PERT~: (min + 4 × plus
          probable + max) / 6.
    \item \textbf{La perte annuelle attendue} compare des options~; elle ne
          décrit pas la pire année.
    \item \textbf{Pertes évitées nettes sur coût}~: le langage du conseil.
\end{itemize}
\end{boxresume}

\begin{boxinfo}
\textbf{Points de contrôle de l'élève.} Je sais~: décomposer un risque
selon FAIR~; nommer les six formes de perte~; calculer une moyenne PERT, une
LEF, une ampleur et une perte annuelle attendue~; comparer deux options en
pertes évitées nettes.
\end{boxinfo}

\bigskip

\begin{boxexo}[title={\ding{45}\quad Exercice~9.1. PERT \niveaubase}]
Un expert estime une perte entre 10 et 70~millions, le plus probablement
25. Calculez la moyenne PERT.
\end{boxexo}

\begin{boxexo}[title={\ding{45}\quad Exercice~9.2. LEF \niveaubase}]
Une fraude au virement est tentée en moyenne 6 fois par an et réussit
une fois sur cinquante. Quelle est la fréquence des pertes~?
\end{boxexo}

\begin{boxexo}[title={\ding{45}\quad Exercice~9.3. Formes de perte \niveaubase}]
Classez dans les six formes de FAIR~: le salaire des agents qui ressaisissent
les factures~; le client qui ne renouvelle pas son contrat~; l'achat d'un
nouveau serveur.
\end{boxexo}

\begin{boxexo}[title={\ding{45}\quad Exercice~9.4. Ampleur \niveaumoyen}]
Refaites le calcul de l'ampleur moyenne de R-02 si la probabilité des
pertes secondaires passe de 0,8 à 0,5.
\end{boxexo}

\begin{boxexo}[title={\ding{45}\quad Exercice~9.5. La fraude \niveaumoyen}]
Avec la LEF de l'exercice 9.2 et une perte moyenne de 60~millions,
calculez la perte annuelle attendue. Une double validation à 4~millions
par an qui divise la réussite par cinq est-elle justifiée~?
\end{boxexo}

\begin{boxexo}[title={\ding{45}\quad Exercice~9.6. Les bornes \niveaumoyen}]
Calculez la perte annuelle attendue de R-02 sans second facteur en
prenant les minima de tous les facteurs, puis les maxima. Que vous
apprend l'écart~?
\end{boxexo}

\begin{boxexo}[title={\ding{45}\quad Exercice~9.7. Construite \niveaumoyen}]
Pourquoi la concordance entre l'estimation et le coût réel de mars est-elle
« en partie construite »~? Comment l'éviter dans une estimation future~?
\end{boxexo}

\begin{boxexo}[title={\ding{45}\quad Exercice~9.8. Moyenne et pire année \niveaumoyen}]
La perte annuelle attendue de R-02 avec second facteur est de 20~millions.
LiZbiZ respecte-t-elle pour autant l'appétence « moins de 100~millions
neuf ans sur dix »~? Que faut-il pour répondre~?
\end{boxexo}

\begin{boxexo}[title={\ding{45}\quad Exercice~9.9. L'inondation \niveauexpert}]
Conduisez une estimation FAIR simplifiée du risque R-11 (inondation de
Bonabéri)~: facteurs, trois points, moyennes, perte annuelle attendue.
\end{boxexo}

\begin{boxexo}[title={\ding{45}\quad Exercice~9.10. La note au conseil \niveauexpert}]
Rédigez la note d'une page de Mme~MOUAFO au conseil sur R-02~: chiffre,
intervalle, hypothèses décisives, recommandation.
\end{boxexo}

\bigskip

\begin{boxinfo}
\textbf{Pour aller plus loin.}
\begin{itemize}
    \item J.~Freund et J.~Jones, \emph{Measuring and Managing Information
          Risk: A FAIR Approach}~\cite{freund2014fair}.
    \item D.~Hubbard et R.~Seiersen~\cite{hubbard2016cyber}.
\end{itemize}
\end{boxinfo}

\begin{boxinfo}
\textbf{Transition.} La moyenne ne dit rien de la pire année. Le
Chapitre~\ref{chap:simulation} propage les incertitudes par simulation de
Monte-Carlo.
\end{boxinfo}
